\section{Results}

\subsection{Clean Image Verification}

Table~\ref{tab:bin_results} reports the verification accuracy on LFW, CFP-FP, and AgeDB-30\cite{cplfw,cfp,agedb}. The teacher model achieves the highest overall accuracy, as expected from the larger MagFace iResNet-100 backbone\cite{magface}. Among the student variants, V1 performs best on the clean bin protocols. In particular, V1/best obtains the highest mean accuracy of 96.37\%, mainly due to stronger performance on CFP-FP and AgeDB-30. V3/best achieves the highest LFW accuracy among the students, but its lower CFP-FP and AgeDB-30 scores reduce its mean performance.

\begin{table}[h]
\centering
\caption{Clean face verification results on LFW, CFP-FP, and AgeDB-30.}
\label{tab:bin_results}

\resizebox{\columnwidth}{!}{%
\begin{tabular}{lcccc}
\hline
\textbf{Model} & \textbf{LFW} & \textbf{CFP-FP} & \textbf{AgeDB-30} & \textbf{Mean} \\
\hline
MagFace iResNet-100 & 99.78 & 96.27 & 98.30 & 98.12 \\
V1/latest & 99.25 & 93.43 & 95.68 & 96.12 \\
V1/best & 99.20 & \textbf{94.14} & \textbf{95.77} & \textbf{96.37} \\
V3/best & \textbf{99.43} & 92.79 & 95.15 & 95.79 \\
V3/latest & 99.02 & 92.06 & 93.60 & 94.89 \\
V3/SWA & 99.00 & 91.54 & 94.00 & 94.85 \\
\hline
\end{tabular}
}
\end{table}

These results show that the clean KD variant is the most suitable checkpoint for standard clean-face verification. Although V3 is designed for robustness, its clean benchmark accuracy is lower than V1, especially on CFP-FP and AgeDB-30. This suggests a trade-off between clean verification accuracy and robustness-oriented training.

\subsection{IJB Template Verification}

Table~\ref{tab:ijb_results} reports template-based verification results on IJB-B and IJB-C. We focus on TAR@$10^{-4}$ because this operating point is more relevant for security-sensitive deployment than average accuracy alone\cite{ijbb,ijbc}.

\begin{table*}[h]
\centering
\caption{IJB-B and IJB-C template verification results. TAR is reported at FAR=$10^{-4}$.}
\label{tab:ijb_results}
\begin{tabular}{lcccc}
\hline
\textbf{Model} & \textbf{IJB-B AUC} & \textbf{IJB-B TAR@$10^{-4}$} & \textbf{IJB-C AUC} & \textbf{IJB-C TAR@$10^{-4}$} \\
\hline
Teacher: MagFace iResNet-100 & 0.9922 & 93.14 & 0.9960 & 97.64 \\
V1/latest & 0.9912 & \textbf{87.98} & 0.9937 & \textbf{90.65} \\
V1/best & 0.9905 & 86.88 & 0.9929 & 89.50 \\
V2/latest & \textbf{0.9935} & 84.52 & \textbf{0.9949} & 86.85 \\
V3/best & 0.9916 & 84.77 & 0.9929 & 87.64 \\
V3/latest & 0.9917 & 84.78 & 0.9932 & 87.23 \\
V3/SWA & 0.9919 & 85.27 & 0.9930 & 87.77 \\
\hline
\end{tabular}
\end{table*}

The teacher remains the strongest model on both IJB-B and IJB-C. Among the student models, V1/latest gives the best low-FAR performance, reaching 87.98\% TAR@$10^{-4}$ on IJB-B and 90.65\% on IJB-C. This confirms that V1 is the preferred student checkpoint for clean or controlled deployment.

An interesting observation is that V2/latest obtains the highest AUC among the student variants but performs worse at TAR@$10^{-4}$. This indicates that AUC alone does not fully reflect strict low-FAR behavior. Since real access-control systems usually operate at low false-acceptance rates, TAR at strict FAR thresholds is more important for deployment.

\subsection{Occlusion Robustness}

Table~\ref{tab:occlusion_results} reports the TAR@$10^{-3}$ drop under synthetic lower-face masking. A smaller absolute drop indicates better robustness. We include MobileFaceNet (W600K) to provide a direct comparison against the lightweight baseline.

\begin{table}[h]
\centering
\caption{Occlusion robustness: TAR@$10^{-3}$ drop (masked$-$clean) under synthetic lower-face masking. Smaller absolute value is better.}
\label{tab:occlusion_results}
\resizebox{\columnwidth}{!}{%
\begin{tabular}{lccccc}
\hline
\textbf{Model} & \textbf{LFW} & \textbf{CFP-FP} & \textbf{AgeDB} & \textbf{CPLFW} & \textbf{CALFW} \\
\hline
MagFace iResNet-100 & -0.041 & -0.520 & -0.541 & -0.724 & -0.236 \\
MobileFaceNet (W600K) & -0.085 & -0.549 & -0.550 & -0.494 & -0.324 \\
V1/latest & -0.037 & -0.369 & -0.462 & -0.222 & -0.248 \\
V3/SWA & \textbf{-0.019} & \textbf{-0.298} & \textbf{-0.321} & \textbf{-0.148} & \textbf{-0.100} \\
\hline
\end{tabular}
}
\end{table}

The key finding is that both student variants outperform MobileFaceNet (W600K) on all five occlusion benchmarks, despite MBF being trained on $7\times$ more identities. MBF suffers the largest TAR drops in this group: $-0.549$ on CFP-FP, $-0.550$ on AgeDB, and $-0.494$ on CPLFW. This is expected because MBF is trained exclusively on clean faces with no occlusion augmentation, leaving it brittle under lower-face masking. In contrast, both V1/latest ($-0.369$, $-0.462$, $-0.222$) and V3/SWA ($-0.298$, $-0.321$, $-0.148$) show substantially smaller degradation, with V3/SWA being the most robust across all protocols.

V3/SWA reduces the absolute TAR drop relative to MBF by $45.7\%$ on CFP-FP, $41.6\%$ on AgeDB, $70.1\%$ on CPLFW, and $69.1\%$ on CALFW. These are the clearest advantages of the proposed occlusion-curriculum and SWA training over a strong clean-face baseline.

The teacher also suffers large degradation under lower-face masking, especially on CFP-FP, AgeDB, and CPLFW. This shows that neither model scale nor training data volume guarantees robustness to structured occlusion\cite{cplfw,cfp,agedb,izmailov2019averagingweightsleadswider}.

\subsection{Overall Comparison}

Table~\ref{tab:checkpoint_recommendation} summarizes the practical role of each checkpoint. V1/latest is the best choice for clean deployment because it gives the strongest student performance on IJB-B and IJB-C at strict FAR. V1/best is preferable when the goal is maximum bin-protocol accuracy. V3/SWA is the best choice when lower-face occlusion or mask usage is expected. V2/latest is kept as an ablation rather than a recommended deployment checkpoint.

\begin{table*}[h]
\centering
\caption{Recommended use cases for the evaluated checkpoints.}
\label{tab:checkpoint_recommendation}
\begin{tabular}{lll}
\hline
\textbf{Checkpoint} & \textbf{Strength} & \textbf{Recommended Use} \\
\hline
V1/latest & Best student IJB TAR@$10^{-4}$ & Clean access-control deployment \\
V1/best & Best bin-protocol mean accuracy & Clean image-pair verification \\
V2/latest & Diagnostic spatial-KD ablation & Not recommended for deployment \\
V3/latest & Robust variant before averaging & Reference checkpoint \\
V3/SWA & Best occlusion robustness & Masked or partially occluded deployment \\
\hline
\end{tabular}
\end{table*}

Overall, the results reveal a clear trade-off. The clean KD variant achieves the best performance under clean verification protocols, while the SWA-stabilized robust variant provides the strongest resistance to lower-face occlusion. Therefore, the appropriate checkpoint should be selected based on the expected deployment environment rather than using a single model for all scenarios.

\subsection{Computational Efficiency}

Table~\ref{tab:efficiency} summarizes the computational characteristics of the teacher, student, and MobileFaceNet baseline. All measurements were taken on a single Tesla P100 16~GB at batch~1 for latency and batch~64 for throughput.

\begin{table}[h]
\centering
\caption{Computational efficiency comparison. Latency and throughput measured on Tesla P100.}
\label{tab:efficiency}
\resizebox{\columnwidth}{!}{%
\begin{tabular}{lcccc}
\hline
\textbf{Model} & \textbf{Params} & \textbf{MACs} & \textbf{Latency (ms)} & \textbf{Memory} \\
\hline
MagFace iResNet-100 & 65.2M & 12.15G & 18.52 & 249 MB \\
\textbf{MobileNetV4-M (student)} & \textbf{9.58M} & \textbf{228M} & \textbf{11.22} & \textbf{38.3 MB} \\
MobileFaceNet (W600K) & $\sim$1M & $\sim$224M & — & — \\
\hline
\end{tabular}
}
\end{table}

The student achieves a 6.8$\times$ parameter reduction and 53.3$\times$ MACs reduction relative to the teacher, while also being faster in wall-clock latency (11.22~ms vs.~18.52~ms at batch~1). The memory footprint is reduced from 249~MB to 38.3~MB, making the student suitable for edge deployment on resource-constrained devices.

\subsection{Lightweight Baseline Comparison}

To contextualize the student's efficiency, we compare against MobileFaceNet trained on WebFace600K (\textasciitilde600K identities) using the insightface \texttt{buffalo\_sc} model with ArcFace loss. Our student uses MS1M-RetinaFace (\textasciitilde85K identities) — 7$\times$ fewer training identities.

\begin{table}[h]
\centering
\caption{Lightweight model comparison on IJB-B/C template verification at FAR=$10^{-4}$.}
\label{tab:baseline_comparison}
\begin{tabular}{lccc}
\hline
\textbf{Model} & \textbf{Params} & \textbf{IJB-B TAR@$10^{-4}$} & \textbf{IJB-C TAR@$10^{-4}$} \\
\hline
Teacher: iResNet-100 & 65.2M & 93.14 & 97.64 \\
MobileFaceNet (W600K) & $\sim$1M & 89.42 & 91.72 \\
Student: V1/latest & 9.58M & 87.98 & 90.65 \\
Student: V3/SWA & 9.58M & 85.27 & 87.77 \\
\hline
\end{tabular}
\end{table}

MobileFaceNet (W600K) has 1M parameters and 7$\times$ more training identities, yet our 9.58M-parameter student is within 1.44\% on IJB-B (87.98 vs.~89.42) and within 1.07\% on IJB-C (90.65 vs.~91.72). Despite having 9$\times$ more parameters than MBF, the student compensates through knowledge distillation from a strong teacher, achieving competitive IJB performance with substantially fewer training identities. The key advantage of the student over MBF lies in occlusion robustness — a regime where MBF, trained only on clean faces, is significantly weaker (see RMFRD evaluation below).

\subsection{Real-World Masked Face Evaluation (RMFRD)}

To evaluate robustness on real masked faces beyond synthetic lower-face masking, we use the RMFRD/AFDB dataset: 403 paired identities with clean gallery images and real-world masked probe images. Images are pre-processed using InsightFace's RetinaFace (\texttt{det\_500m}) detector to obtain 5-point face landmarks and apply standard ArcFace affine alignment to 112$\times$112.

\textbf{Protocol:} Gallery = mean embedding of up to 10 clean images per identity; Probes = all masked images; Positive pairs = same-identity mask-to-clean matching; TAR@FAR and Rank-1 identification are reported.

\begin{table}[h]
\centering
\caption{Real-world masked face recognition on RMFRD (403 paired identities, RetinaFace-aligned). Gallery: mean of up to 10 clean images per identity. Probe: 1,945 real masked images. Fallback bilinear resize for 7.9\% of masked images where RetinaFace detected no face. Unaligned reference: MBF 68.72\% AUC / 7.25\% Rank-1; V1/best 71.02\% / 4.78\%; V3/SWA 69.90\% / 6.84\%.}
\label{tab:rmfrd}
\resizebox{\columnwidth}{!}{%
\begin{tabular}{lccc}
\hline
\textbf{Model} & \textbf{AUC} & \textbf{TAR@$10^{-3}$} & \textbf{Rank-1} \\
\hline
MobileFaceNet (W600K) & 85.55 & 26.43 & 36.56 \\
Student: V1/best & 82.66 & 15.06 & 22.01 \\
Student: V3/SWA & \textbf{84.18} & \textbf{15.68} & \textbf{23.03} \\
\hline
\end{tabular}
}
\end{table}

The alignment step is critical for meaningful evaluation. Without alignment, all models use simple bilinear resize, which suppresses discriminative facial features: MBF achieves only 68.72\% AUC / 7.25\% Rank-1 on unaligned crops but reaches 85.55\% AUC / 36.56\% Rank-1 after RetinaFace alignment — an improvement of $+16.8\%$ AUC and $+29.3$ Rank-1 points. This confirms that the RMFRD raw images are not pre-aligned and that any meaningful comparison must use landmark-based preprocessing.

Among the student models, V3/SWA (84.18\% AUC, 23.03\% Rank-1) outperforms V1/best (82.66\% AUC, 22.01\% Rank-1) on real masked probes, confirming that the occlusion curriculum and SWA stabilization generalize beyond synthetic lower-face masking to real-world masks. This RMFRD result directly complements the synthetic masking ablation (Table~\ref{tab:occlusion_results}): just as V3/SWA shows smaller TAR drop on synthetic occlusion, it also achieves higher AUC and Rank-1 on real masked faces.

MobileFaceNet (W600K), despite its smaller parameter count (1M vs.\ 9.58M), outperforms both student variants on this dataset (85.55\% AUC, 36.56\% Rank-1). We attribute this primarily to the significantly larger training set: MBF is trained on WebFace600K with $\sim$600K identities, while our student uses MS1M-RetinaFace with $\sim$85K identities (7$\times$ fewer). Additionally, the RMFRD dataset predominantly features Asian celebrities, who may be more represented in WebFace600K than in MS1M. This suggests that the student's IJB advantage does not fully translate to this particular domain, highlighting the importance of training data diversity for cross-domain robustness.

\subsection{Real-World Masked Face Evaluation (MFR2)}

MFR2 is a real-world masked face dataset of 55 celebrity identities with both masked and unmasked images per subject. We evaluate 53 identities that have at least one clean (\texttt{no-mask}) and one masked image. Two protocols are evaluated:

\begin{itemize}
    \item \textbf{Verification (pairs.txt)}: 424 same-identity and 424 different-identity pairs sourced from the official LFW-style split; pairs may combine masked and unmasked images. We report AUC.
    \item \textbf{1:N Identification}: gallery = mean embedding of all clean images per identity (98 total); probes = all masked images (171 total). We report AUC and Rank-1.
\end{itemize}

Note: With only 424 negative pairs in the verification split, the FAR resolution ($\approx 1/424 \approx 0.0024$) is too coarse to report meaningful TAR at FAR$\leq10^{-3}$; we therefore report AUC only for MFR2.

\begin{table}[h]
\centering
\caption{MFR2 evaluation (53 paired identities). Verification: pairs.txt AUC (mixed masked/clean pairs). Identification: 1:N AUC and Rank-1 with clean gallery and masked probes.}
\label{tab:mfr2}
\resizebox{\columnwidth}{!}{%
\begin{tabular}{lccc}
\hline
\textbf{Model} & \textbf{Verif AUC} & \textbf{ID AUC} & \textbf{ID Rank-1} \\
\hline
MobileFaceNet (W600K) & 89.91 & 95.11 & \textbf{72.51} \\
Student: V1/best      & 94.18 & 95.37 & 61.99 \\
Student: V3/SWA       & \textbf{94.38} & \textbf{95.99} & 66.67 \\
\hline
\end{tabular}
}
\end{table}

Both student models substantially outperform MobileFaceNet (W600K) on MFR2 verification AUC (V1/best: $+4.27\%$; V3/SWA: $+4.47\%$). V3/SWA also surpasses MBF on identification AUC ($+0.88\%$). The improvement on verification AUC confirms that the student's knowledge-distilled embeddings yield stronger pairwise discriminability on a mixed masked/unmasked face matching task, despite training on $7\times$ fewer identities.

However, MBF retains a substantial lead on Rank-1 identification ($72.51\%$ vs.\ $66.67\%$ for V3/SWA). This mirrors the RMFRD pattern, where MBF also leads on Rank-1. We attribute MBF's Rank-1 advantage to its larger training set ($\sim\!600$K identities vs.\ $\sim\!85$K) and, potentially, to greater overlap with the post-2020 political celebrity identities featured in both MFR2 and RMFRD. The student's superior AUC suggests better relative embedding calibration, but fewer absolute successes at the top-1 threshold.

\subsection{Runtime Pipeline Results}

The runtime pipeline was designed to evaluate the trained encoder under realistic video conditions. It records processed frames, detection counts, accepted observations, recognized observations, match rejection reasons, and mean FPS. However, since no standardized annotated video benchmark is used in the current evaluation, we do not report a single comparable runtime accuracy table in this section. Instead, the benchmark results above focus on the recognition quality of the trained encoders, while the runtime pipeline is used to validate practical deployment behavior with detection, tracking, liveness filtering, quality gating, pooling, and FAISS retrieval\cite{faiss}.